\documentclass[../main.tex]{subfiles}
\begin{document}
\section{Introduction}
\label{sec:introduction}
\label{subs:background}
The credibility of empirical research in econometrics is fundamentally tied to the robustness of its findings. However, recent developments have shown that statistical results are frequently driven by incredibly small subsets of data, making them highly sensitive to specific observations \parencite{broderick2020automatic}. While researchers typically rely on some classical diagnostic tools or robust estimators to maintain integrity, both approaches have limitations. Diagnostic tools often fail to capture complex higher-order interactions such as joint influence and masking \parencite{chatterjee1986influential, rousseeuw1990unmasking}. Conversely, robust estimators prioritize controlling worst-case sensitivity to stabilize the overall model, rather than identifying the specific and minimal subsets of data that actually drive a particular statistical conclusion. \parencite{hampel1974influence}.
This study addresses this gap by utilizing the exact Most Influential Sets (MIS) tool to redefine how we assess empirical robustness \parencite{kuschnig2021hidden}.

\label{subs:problem_statement}
Traditionally, leave-one-out measures such as Cook's distance and DFBETAs quantify influence by observing the standardized change in parameter vectors when a single observation is deleted. While valuable as local sensitivity metrics, this single-observation design leaves them susceptible to "masking" effects, where influential observations mutually conceal each other's impact \parencite{cook2000detection}. Similarly, leverage metrics derived from the diagonal elements of the hat matrix ($h_{ii}$) assess how mathematically extreme an observation is within the multidimensional predictor space. However, these tools possess a breakdown point of $1/n$, meaning a single sufficiently extreme outlier or a cluster of outliers can render the diagnostic entirely unreliable \parencite{rousseeuw2003robust}. This failure is primarily driven by the masking effect: because classical methods rely on the sample mean and non-robust covariance matrices, a cluster of outliers can distort these foundational moments \parencite{maronna2019robust}. Consequently, the presence of a set of anomalies artificially deflates both the leverage scores and the deletion metrics of individual points within that cluster, allowing structurally influential sets to remain hidden from standard diagnostic audits \parencite{rousseeuw1990unmasking}.

\label{subs:objectives}

\begin{enumerate}
	\item \textbf{Detection of Higher-Order Interactions:}
	\item \textbf{Algorithmic Validation and EVT Integration:}
	\item \textbf{Empirical Resolution of Controversies:}
\end{enumerate}


\label{subs:method_summary}


\label{subs:thesis sturcture}

\end{document}